\section{Introduction}
La préparation d'un devis de travaux demande de comprendre une demande parfois imprécise, de sélectionner les opérations pertinentes, de renseigner les variantes techniques et d'établir un prix cohérent. Cette activité combine vocabulaire métier, contraintes de catalogue et calculs de coûts. Lorsqu'elle est réalisée manuellement, elle mobilise du temps et peut produire des écarts entre le besoin exprimé par le client et les lignes retenues dans le devis.

Le projet \textit{Devis IA} répond à cette problématique par un système multi-modal de génération assistée de devis. Le commercial peut décrire les travaux au clavier, dicter sa demande ou fournir une photographie. Le système produit ensuite des occurrences de prestations reliées au catalogue réel de l'entreprise, puis prépare un devis contrôlable par un utilisateur métier.

L'enjeu n'est pas seulement de produire un texte plausible. La sortie de l'intelligence artificielle doit respecter une structure JSON stricte, utiliser des identifiants réels de la base de données et préserver les relations entre prestations, options et choix. Le reste de ce rapport décrit le besoin, l'architecture, les modalités développées, le protocole d'évaluation et les résultats actuellement disponibles.
